\documentclass[../../../include/open-logic-section]{subfiles}

\begin{document}

\olfileid{sth}{z}{infinity-again}
\olsection{لامتناہیت}

ساڈے کول پہلاں ای اینے مسلّمے نیں کہ جے کوئی سیٹ موجود اے، تاں لامتناہی
گنتی وچ سیٹ ہون گے۔ کیوں جے منّ لو کوئی سیٹ موجود اے، تاں
$\emptyset$ موجود اے (\olref[sth][z][sep]{prop:emptyexists} نال)۔
ہُن کسے وی سیٹ~$x$ لئی سیٹ $x \cup \{x\}$،
\olref[sth][z][pairs]{prop:pairsconsequences} نال موجود اے۔
سو ایہ گل کجھ واری لا کے سانوں ایہ سیٹ ملن گے:
\begin{enumerate}
	\item[0.] $\emptyset$ %& $0$ & stage $0$\\
	\item[1.] $ \{\emptyset\}$ %& $1$ & stage $1$\\
	\item[2.] $\{\emptyset, \{\emptyset\}\}$ %& $2$ & stage $2$\\
	\item[3.] $\{\emptyset, \{\emptyset\}, \{\emptyset, \{\emptyset\}\}\}$ %&$3$ & stage $3$\\
	\item[4.] $\{\emptyset, \{\emptyset\}, \{\emptyset, \{\emptyset\}\}, \{\emptyset, \{\emptyset\}, \{\emptyset, \{\emptyset\}\}\}\}$% & $4$ & stage $4$ 
\end{enumerate}
تے اسیں پرکھ سکدے آں کہ ایہناں وچوں ہر سیٹ وکھرا اے۔

اسیں گنتی $0$ توں کئی وجوہاں کر کے شروع کیتی اے۔ اوہناں وچوں اک ایہ اے۔
ایہ پرکھنا بہتا اوکھا نہیں کہ جیہڑے سیٹ نوں اسیں ’’$n$‘‘ دا نشان دتا اے،
اوہدے عین $n$ رکن نیں، تے سمجھ مطابق اوہ $n$ویں مرحلے اُتے بندا اے۔

پر ایس نال سانوں \emph{لامتناہی گنتی وچ} سیٹ ملدے نیں۔ ایہ ضمانت نہیں
ملدی کہ اک \emph{لامتناہی سیٹ} وی اے، یعنی اک سیٹ جیہدے رکن لامتناہی
گنتی وچ ہون۔ تے ایس نال واقعی فرق پیندا اے۔ جِنّی دیر سانوں اک (ڈیڈیکنڈ)
لامتناہی سیٹ نہیں لبھدا، اسیں ڈیڈیکنڈ الجبرا نہیں بنا سکدے۔ پر سانوں
ڈیڈیکنڈ الجبرا چاہیدا اے تاں جو اسیں اوہنوں قدرتی عدداں دا سیٹ منّ سکیے۔
(\olref[sfr][infinite][dedekindsproof]{sec} نال ملا کے ویکھو۔)

خاص گل ایہ اے کہ ہُن تک دتے مسلّمے کسے وی لامتناہی سیٹ دی موجودگی دی
ضمانت \emph{نہیں} دیندے۔ سو سانوں اک نواں مسلّمہ دینا پئے گا:

\begin{axiom}[لامتناہیت]
اک سیٹ $I$ موجود اے جیہدے لئی $\emptyset \in I$ تے $x \cup \{x\} \in I$ ہر اوس صورت وچ جتھے $x \in I$۔
\begin{align*}
	\exists I( & (\exists o \in I)\forall x\ x \notin o \land {}\\
	& (\forall x \in I)(\exists s \in I)\forall z(z \in s \liff (z \in x \lor z = x)))
\end{align*}
\end{axiom}

ایہ ویکھنا سوکھا اے کہ لامتناہیت دے مسلّمے توں ملیا سیٹ $I$ ڈیڈیکنڈ
لامتناہی اے۔ اوہدا خاص رکن $\emptyset$ اے تے $I$ اُتے یک بہ یک نقشہ
$s(x) = x\cup \{x\}$ اے۔ ہُن \olref[sfr][infinite][dedekind]{thm:DedekindInfiniteAlgebra}
نے وکھایا کہ ڈیڈیکنڈ لامتناہی سیٹ توں ڈیڈیکنڈ الجبرا کیویں ملدا اے۔ اسیں
ایہنوں قدرتی عدداں دا اپنا سیٹ منّاں گے۔ ہور ٹھیک لفظاں وچ:
\begin{defn}\ollabel{defnomega}
منّ لو $I$ لامتناہیت دے مسلّمے توں ملیا کوئی وی سیٹ اے۔ منّ لو $s$
فنکشن $s(x) = x \cup \{x\}$ اے۔ منّ لو $\omega = \closureofunder{s}{\emptyset}$۔
اسیں $\omega$ دے رکناں نوں \emph{قدرتی عدد} آکھدے آں، تے آکھدے آں کہ $n$
اوہ نتیجہ اے جیہڑا $n$ واری $s$ نوں $\emptyset$ اُتے لاؤن نال ملدا اے۔
\end{defn}

ہُن تسی پچھاں ویکھ کے پرکھ سکدے او کہ چند پیرا پہلاں ’’$n$‘‘ دا نشان
دتا سیٹ، عدد $n$ \emph{دے طور اُتے} ورتیا جائے گا۔

اسیں ایس مقرر کیتی گل دی اہمیت \olref[sth][z][nat]{sec} وچ ویکھاں گے۔
ہُن ایس نال اسیں سمجھ آؤن والا اک نتیجہ ثابت کر سکدے آں:

\begin{prop}\ollabel{naturalnumbersarentinfinite}
کوئی قدرتی عدد ڈیڈیکنڈ لامتناہی نہیں۔
\end{prop}

\begin{proof}
ثبوت استقرا نال اے، یعنی \olref[sfr][infinite][induction]{thm:dedinfiniteinduction} نال۔
صاف اے کہ $0 = \emptyset$ ڈیڈیکنڈ لامتناہی نہیں۔ استقرائی قدم لئی اسیں
عکسِ نقیض ثابت کراں گے۔ جے (بے تُکے مفروضے مطابق) $s(n)$ ڈیڈیکنڈ لامتناہی
اے، تاں $n$ ڈیڈیکنڈ لامتناہی اے۔

سو منّ لو $s(n)$ ڈیڈیکنڈ لامتناہی اے۔ یعنی اک !!{injection} $f$ اے
جیہدے لئی $\ran{f}\subsetneq \dom{f} = s(n) = n \cup \{n\}$۔
دونیں صورتاں ویکھدے آں۔

\emph{پہلی صورت۔ $n \notin \ran{f}$۔} سو $\ran{f} \subseteq n$ تے
$f(n) \in n$۔ منّ لو $g = \funrestrictionto{f}{n}$۔ ہُن
$\ran{g} = \ran{f} \setminus \{f(n)\} \subsetneq n = \dom{g}$۔
ایس لئی $n$ ڈیڈیکنڈ لامتناہی اے۔

\emph{دوجی صورت۔ $n \in \ran{f}$۔} $m \in \dom{f} \setminus \ran{f}$ مقرر
کرو تے فنکشن $h$ دی تعریف دائرہ $s(n) = n \cup \{n\}$ اُتے کرو:
\[
h(x) =
	\begin{cases}
		f(x) & \text{جے }f(x) \neq n\\
		m & \text{جے }f(x)=n
	\end{cases}
\]
سو $h$ تے $f$ ہر تھاں ملدے نیں، بس $h(f^{-1}(n)) = m \neq n = f(f^{-1}(n))$۔
چونکہ $f$ اک !!a{injection} اے، ایس لئی $n \notin \ran{h}$ تے
$\ran{h} \subsetneq \dom{h} = s(n)$۔ ہُن پہلی صورت والی دلیل نال
$n$ ڈیڈیکنڈ لامتناہی اے۔
\end{proof}

پر سوال رہ جاندا اے کہ اسیں لامتناہیت دے مسلّمے دا \emph{جواز} کیویں دیئیے۔
مختصر جواب ایہ اے کہ ساڈے ہُن تک دے بیان وچ اک ہور اصول شامل کرنا پئے گا۔
اصول ایہ اے:
\begin{enumerate}
	\item[] \stagesinf۔ اک لامتناہی مرحلہ اے۔ یعنی اک ایہو جیہا مرحلہ اے جیہڑا (ا) پہلا مرحلہ نہیں، (ب) جیہدے توں پہلاں کجھ مرحلے نیں، پر (ج) جیہدے توں عین پہلاں کوئی مرحلہ نہیں۔
\end{enumerate}
لامتناہیت دا مسلّمہ ایس اصول توں سدھا نکلدا اے۔ سانوں پتا اے کہ قدرتی عدد
$n$ مرحلے $n$ اُتے بندا اے۔ سو سیٹ $\omega$ پہلے لامتناہی مرحلے اُتے
بندا اے۔ تے $\omega$ آپ لامتناہیت دے مسلّمے دا گواہ اے۔

پر ایس نال اسیں واپس ایس سوال ول آ جاندے آں کہ \stagesinf دا جواز
کیویں دیئیے۔ \stagessucc وانگ ایہ ساڈے اکٹھا کرن والے مرحلہ وار درجیاں
دے بیان دا صاف حصہ نہیں سی۔ ایس توں وی ودھ گل ایہ اے کہ مرحلہ وار درجیاں
دے خیال وچ، جیہدے وچ سیٹ واری واری مرحلیاں اُتے بندے نیں، کوئی گل سانوں
ایہ منّن لئی مجبور نہیں کردی کہ کوئی \emph{لامتناہی} مرحلہ اے۔ ایہ سوچنا
بالکل ٹھیک بیٹھدا اے کہ مرحلے قدرتی عدداں وانگ ترتیب وچ نیں۔

پر ایس نال اک صاف مسئلہ پیدا ہوندا اے۔ \olref[sfr][infinite][dedekindsproof]{sec}
وچ اسیں ڈیڈیکنڈ دا ’’ثبوت‘‘ ویکھیا سی کہ (خیالاں دا) اک ڈیڈیکنڈ لامتناہی
سیٹ اے۔ ہو سکدا اے ایہ تہانوں تسلی بخش نہ لگیا ہووے۔ پر جے سیٹ دا مرحلہ
وار تصور (یا ’’سوچ دے قانون‘‘) سانوں \stagesinf{} منّن لئی ’’مجبور نہیں
کردے‘‘، تاں ڈیڈیکنڈ لامتناہی سیٹ دی موجودگی دے دعوے دا اپنا اندرونی جواز
ساڈے کول ہجے وی نہیں۔

بے شک ایتھے کہن نوں ہور بہت کجھ اے۔ پر آس اے کہ ہُن تسی سوچنا شروع کر
سکدے او کہ کسے مسلّمے (یا اصول) دا جواز دین لئی کیہدی \emph{لوڑ} پَندی اے۔
اگے اسیں \stagesinf{} نوں بس منّ لواں گے۔

\end{document}
